\CNsection{1}{User Plane Concept}

The \textbf{User Plane} (also called the \textbf{Data Plane} or \textbf{U-Plane}) is responsible for carrying actual user data through the mobile network. This includes:

\begin{itemize}
\tightlist
\item
  \textbf{IP packets} (web browsing, app data, file transfers)
\item
  \textbf{Voice data} (VoLTE/\allowbreak{}VoNR encoded speech frames)
\item
  \textbf{Video streams} (real-time video, streaming content)
\item
  \textbf{Any application-layer payload}
\end{itemize}

The user plane is distinct from the control plane — while the control plane sets up, manages, and tears down connections, the user plane is the ``pipe'' through which data actually flows once those connections are established.

\textbf{Key characteristics:}

\begin{itemize}
\tightlist
\item
  High throughput requirement
\item
  Low latency sensitivity
\item
  Stateless forwarding (once tunnels are established)
\item
  QoS-aware packet handling
\item
  Transparent to signaling — carries data without interpreting it
\end{itemize}

\CNkeepfig{m23-gtpu-header}{12}

\CNsection{2}{GTP-U Protocol (GPRS Tunneling Protocol — User Plane)}

GTP-U (3GPP TS 29.281) is the core tunneling protocol used to transport user data between network nodes across all generations (3G, 4G, 5G).

\CNsubsection{2.1}{Header Structure}

The GTP-U header has an \textbf{8-byte mandatory portion} plus optional extension headers:

\CNfigure{m23-gtpu-header}{GTP-U header}

{\def\LTcaptype{none} % do not increment counter
\Needspace{20\baselineskip}
\begin{longtable}[c]{@{}ccc@{}}
\toprule\noalign{}
\rowcolor{CNink}\CNth{Field} & \CNth{Bits} & \CNth{Description} \\
\CNheadrulesep
\endhead
\bottomrule\noalign{}
\endlastfoot
Version & 3 & Always \CNcode{001} (GTPv1) \\
PT (Protocol Type) & 1 & \CNcode{1} = GTP, \CNcode{0} = GTP’ (charging) \\
Reserved (*) & 1 & Must be 0 \\
E (Extension Header) & 1 & Extension header follows \\
S (Sequence Number) & 1 & Sequence number field present \\
PN (N-PDU Number) & 1 & N-PDU number field present \\
Message Type & 8 & \CNcode{0xFF} = T-PDU (user data), others = signaling \\
Length & 16 & Payload length (excludes mandatory 8~bytes) \\
TEID & 32 & Tunnel Endpoint Identifier \\
\end{longtable}
}

\CNsubsection{2.2}{TEID (Tunnel Endpoint Identifier)}

The \textbf{TEID} is a 32-bit value that uniquely identifies a GTP tunnel endpoint at a receiving node:

\begin{itemize}
\tightlist
\item
  Each direction of a tunnel has its own TEID
\item
  The \textbf{receiving node} assigns the TEID (tells the sender ``use this TEID when sending to me'')
\item
  TEID = 0 is reserved for control plane initial messages
\item
  A node may have thousands of TEIDs active simultaneously (one per bearer/\allowbreak{}session per direction)
\end{itemize}

\textbf{Example:}

\tcbset{CNcodemode/.style={unbreakable}}\def\CNcodelang{}

\begin{CNplaincode}
\begin{Verbatim}
gNB tells UPF: "For downlink data to UE-X, use TEID = 0x00A1B2C3 on N3"
UPF tells gNB: "For uplink data from UE-X, use TEID = 0x00D4E5F6 on N3"
\end{Verbatim}
\end{CNplaincode}

\CNkeepfig{m23-gtpu-encap}{5}

\CNsubsection{2.3}{Encapsulation /\allowbreak{} Decapsulation}

At each GTP-U hop, the user’s IP packet is \textbf{encapsulated} inside:

\CNfigure{m23-gtpu-encap}{GTP-U encapsulation at each hop}

\begin{itemize}
\tightlist
\item
  \textbf{Encapsulation} (sending node): Prepends GTP-U header + UDP + outer IP
\item
  \textbf{Decapsulation} (receiving node): Strips outer IP + UDP + GTP-U header, forwards inner IP
\item
  At intermediate nodes: may decapsulate and re-encapsulate with new TEID
\end{itemize}

\CNsubsection{2.4}{Sequence Numbers and Extension Headers}

\textbf{Sequence Numbers (16-bit):}

\begin{itemize}
\tightlist
\item
  Optional but used for reordering at tunnel endpoints
\item
  Critical during handover (ensures no packet loss/\allowbreak{}reorder)
\item
  Incremented per PDU per tunnel
\end{itemize}

\textbf{Extension Headers:}

\begin{itemize}
\tightlist
\item
  Chained via ``Next Extension Header Type'' field
\item
  Each extension: \CNcode{[Length (1 byte)][Content (variable)][Next Ext Hdr Type (1 byte)]}
\end{itemize}

\Needspace{12\baselineskip}

Common extension headers:

{\def\LTcaptype{none} % do not increment counter
\Needspace{9\baselineskip}
\begin{longtable}[c]{@{}
  >{\raggedright\arraybackslash\hspace{0pt}}m{(\linewidth - 4\tabcolsep) * \real{0.2857}}
  >{\raggedright\arraybackslash\hspace{0pt}}m{(\linewidth - 4\tabcolsep) * \real{0.2857}}
  >{\raggedright\arraybackslash\hspace{0pt}}m{(\linewidth - 4\tabcolsep) * \real{0.4286}}@{}}
\toprule\noalign{}
\rowcolor{CNink}\CNth{Type} & \CNth{Name} & \CNth{Purpose} \\
\CNheadrulesep
\endhead
\bottomrule\noalign{}
\endlastfoot
0x85 & PDU Session Container & 5G QoS Flow Identifier (QFI), reflective QoS \\
0xC0 & PDCP PDU Number & For lossless handover \\
0x40 & UDP Port & Additional multiplexing \\
\end{longtable}
}

\Needspace{14\baselineskip}

\CNkeepfig{m23-lte-path}{8}

\CNsection{3}{User Plane in LTE (4G)}

\CNchained\CNsubsection{3.1}{End-to-End Path}

\CNfigure{m23-lte-path}{End-to-end LTE user-plane path}

\CNkeepfig{m23-lte-segments}{3}

\CNsubsection{3.2}{Packet Headers at Each Segment}

\CNfigure{m23-lte-segments}{Packet headers on each LTE user-plane segment}

\CNsubsection{3.3}{LTE Bearer Model}

\begin{itemize}
\tightlist
\item
  \textbf{Default Bearer}: Always-on, established at attach (one per PDN connection)
\item
  \textbf{Dedicated Bearer}: Additional QoS, established on demand (e.g., VoLTE)
\item
  Each bearer = unique set of TEIDs on S1-U and S5
\end{itemize}

\Needspace{14\baselineskip}

\CNkeepfig{m23-5g-path}{8}

\CNsection{4}{User Plane in 5G}

\CNchained\CNsubsection{4.1}{End-to-End Path}

\CNfigure{m23-5g-path}{End-to-end 5G user-plane path}

\Needspace{13\baselineskip}

\CNsubsection{4.2}{5G Reference Points for User Plane}

{\def\LTcaptype{none} % do not increment counter
\Needspace{9\baselineskip}
\begin{longtable}[c]{@{}cccc@{}}
\toprule\noalign{}
\rowcolor{CNink}\CNth{Interface} & \CNth{Between} & \CNth{Transport} & \CNth{Purpose} \\
\CNheadrulesep
\endhead
\bottomrule\noalign{}
\endlastfoot
\textbf{N3} & gNB \ensuremath{\leftrightarrow} UPF & GTP-U over UDP/\allowbreak{}IP & RAN to Core user data \\
\textbf{N9} & UPF \ensuremath{\leftrightarrow} UPF & GTP-U over UDP/\allowbreak{}IP & Inter-UPF forwarding \\
\textbf{N6} & UPF \ensuremath{\leftrightarrow} DN & Native IP & Exit to data network \\
\end{longtable}
}

\CNkeepfig{m23-nr-stack}{5}

\CNsubsection{4.3}{SDAP Layer (Service Data Adaptation Protocol)}

The \textbf{SDAP} layer is new in 5G NR (not present in LTE) and sits above PDCP:

\CNfigure{m23-nr-stack}{NR user-plane stack with the new SDAP layer}

\textbf{SDAP Functions:}

\begin{itemize}
\tightlist
\item
  \textbf{Downlink:} QoS Flow \ensuremath{\to} DRB mapping (based on QFI in GTP-U extension header)
\item
  \textbf{Uplink:} DRB \ensuremath{\to} QoS Flow marking (stamps QFI on uplink packets)
\item
  \textbf{Reflective QoS:} UE mirrors the DL mapping for UL without explicit signaling
\end{itemize}

\Needspace{14\baselineskip}

\CNsubsection{4.4}{Multiple UPFs Architecture}

5G supports \textbf{UPF chaining} for flexible traffic routing:

{\def\LTcaptype{none} % do not increment counter
\Needspace{8\baselineskip}
\begin{longtable}[c]{@{}
  >{\raggedright\arraybackslash\hspace{0pt}}m{(\linewidth - 4\tabcolsep) * \real{0.3846}}
  >{\raggedright\arraybackslash\hspace{0pt}}m{(\linewidth - 4\tabcolsep) * \real{0.2308}}
  >{\raggedright\arraybackslash\hspace{0pt}}m{(\linewidth - 4\tabcolsep) * \real{0.3846}}@{}}
\toprule\noalign{}
\rowcolor{CNink}\CNth{UPF Type} & \CNth{Role} & \CNth{Use Case} \\
\CNheadrulesep
\endhead
\bottomrule\noalign{}
\endlastfoot
\textbf{I-UPF} (Intermediate) & Transit/\allowbreak{}branching point & Handover anchor, traffic steering \\
\textbf{PSA-UPF} (PDU Session Anchor) & Session anchor, IP allocation & Stable IP address point \\
\end{longtable}
}

\CNfigure{m23-upf-chain}{UPF chaining with a local breakout to the edge}

\Needspace{24\baselineskip}

\CNsection{5}{UPF Functions (3GPP TS 23.501)}

The UPF is controlled by the SMF via the \textbf{N4 interface} using the \textbf{PFCP protocol} (Packet Forwarding Control Protocol, TS 29.244).

\CNsubsection{5.1}{PFCP Rules}

{\def\LTcaptype{none} % do not increment counter
\Needspace{13\baselineskip}
\begin{longtable}[c]{@{}
  >{\raggedright\arraybackslash\hspace{0pt}}m{(\linewidth - 4\tabcolsep) * \real{0.2400}}
  >{\raggedright\arraybackslash\hspace{0pt}}m{(\linewidth - 4\tabcolsep) * \real{0.3600}}
  >{\raggedright\arraybackslash\hspace{0pt}}m{(\linewidth - 4\tabcolsep) * \real{0.4000}}@{}}
\toprule\noalign{}
\rowcolor{CNink}\CNth{Rule} & \CNth{Acronym} & \CNth{Function} \\
\CNheadrulesep
\endhead
\bottomrule\noalign{}
\endlastfoot
\textbf{Packet Detection Rule} & PDR & Match incoming packets (by TEID, IP, port, QFI, etc.) \\
\textbf{Forwarding Action Rule} & FAR & What to do (forward, drop, buffer, duplicate) \\
\textbf{QoS Enforcement Rule} & QER & Rate limiting (MBR/\allowbreak{}GBR), gating, marking \\
\textbf{Usage Reporting Rule} & URR & Volume/\allowbreak{}time measurement, triggers for reporting \\
\textbf{Buffering Action Rule} & BAR & Buffer DL packets when UE is idle, notify SMF \\
\end{longtable}
}

\CNkeepfig{m23-pdr-rules}{3}

\CNsubsection{5.2}{How Rules Work Together}

\CNfigure{m23-pdr-rules}{How the PFCP rules work together in the UPF}

\CNsubsection{5.3}{N4/\allowbreak{}PFCP Interface}

\begin{itemize}
\tightlist
\item
  \textbf{SMF \ensuremath{\to} UPF:} Session Establishment/\allowbreak{}Modification/\allowbreak{}Deletion Request
\item
  \textbf{UPF \ensuremath{\to} SMF:} Session Report (usage reports, DL data notification)
\item
  PFCP uses \textbf{SEID} (Session Endpoint Identifier) — analogous to TEID but for control
\item
  Runs over UDP port 8805
\end{itemize}

\CNsubsection{5.4}{Complete UPF Function List}

\begin{enumerate}
\def\labelenumi{\arabic{enumi}.}
\tightlist
\item
  Packet routing \& forwarding
\item
  Packet inspection (DPI for application detection)
\item
  UL/\allowbreak{}DL rate enforcement (per QoS flow)
\item
  Traffic usage reporting (for charging)
\item
  DL packet buffering \& notification
\item
  QoS handling (DSCP marking, policing)
\item
  Lawful intercept (UP collection)
\item
  Transport level packet marking (outer header DSCP)
\end{enumerate}

\CNkeepfig{m23-dl-steps}{10}

\CNsection{6}{Packet Trace Example: Downlink Data}

\textbf{Scenario:} Server on Internet sends data to UE

\CNsubsection{6.1}{Step-by-Step}

\CNfigure{m23-dl-steps}{A downlink packet at each step from the Internet to the UE}

\CNkeepfig{m23-encapsulation}{3}

\CNsubsection{6.2}{Encapsulation Layer Diagram}

\CNfigure{m23-encapsulation}{Headers of a downlink packet at each hop}

\Needspace{14\baselineskip}

\CNkeepfig{m23-split-bearer}{10}

\CNsection{7}{User Plane Optimization}

\CNchained\CNsubsection{7.1}{Dual Connectivity (DC) — Split Bearer}

In \textbf{EN-DC} (E-UTRA–NR Dual Connectivity) or \textbf{NR-DC}:

\CNfigure{m23-split-bearer}{Dual connectivity with a split bearer}

\Needspace{12\baselineskip}

\textbf{Split Bearer types:}

{\def\LTcaptype{none} % do not increment counter
\Needspace{9\baselineskip}
\begin{longtable}[c]{@{}
  >{\raggedright\arraybackslash\hspace{0pt}}m{(\linewidth - 4\tabcolsep) * \real{0.1935}}
  >{\raggedright\arraybackslash\hspace{0pt}}m{(\linewidth - 4\tabcolsep) * \real{0.4194}}
  >{\raggedright\arraybackslash\hspace{0pt}}m{(\linewidth - 4\tabcolsep) * \real{0.3871}}@{}}
\toprule\noalign{}
\rowcolor{CNink}\CNth{Type} & \CNth{Description} & \CNth{GTP-U path} \\
\CNheadrulesep
\endhead
\bottomrule\noalign{}
\endlastfoot
MCG Bearer & Data only via Master Node & UPF \ensuremath{\leftrightarrow} MN \\
SCG Bearer & Data only via Secondary Node & UPF \ensuremath{\leftrightarrow} SN \\
Split Bearer & Data split across both nodes & UPF \ensuremath{\leftrightarrow} MN \ensuremath{\leftrightarrow} SN (MN routes) \\
\end{longtable}
}

\begin{itemize}
\tightlist
\item
  Increases throughput by aggregating radio resources from two nodes
\item
  PDCP layer at MN handles reordering for split bearers
\end{itemize}

\CNsubsection{7.2}{Multi-Path UPF}

Multiple UPFs in the data path for:

\begin{itemize}
\tightlist
\item
  \textbf{Load balancing}: Distribute traffic across UPFs
\item
  \textbf{Redundancy}: N9 links between UPFs for resiliency
\item
  \textbf{Geographic optimization}: Route traffic through nearest UPF
\end{itemize}

\CNkeepfig{m23-ulcl}{7}

\CNsubsection{7.3}{Edge Computing (Local Breakout)}

5G enables \textbf{traffic steering} for MEC (Multi-access Edge Computing):

\textbf{UL Classifier (UL CL):}

\CNfigure{m23-ulcl}{Uplink classifier for edge breakout}

\begin{itemize}
\tightlist
\item
  UL CL inspects UL packets and routes matching traffic to local DN
\item
  Non-matching traffic continues to central DN
\item
  \textbf{Same UE IP address} for both paths (transparent to UE)
\end{itemize}

\CNkeepfig{m23-branching}{2}

\textbf{Branching Point:}

\CNfigure{m23-branching}{Branching point for IPv6 multi-homing}

\begin{itemize}
\tightlist
\item
  Different PDU session anchors for local vs. remote
\item
  May use different UE IP addresses
\end{itemize}

\textbf{Benefits of Edge Computing:}

\begin{itemize}
\tightlist
\item
  Ultra-low latency (\textless{} 10~ms RTT)
\item
  Reduced backhaul bandwidth
\item
  Local data processing (privacy, compliance)
\item
  Use cases: AR/\allowbreak{}VR, autonomous driving, industrial IoT
\end{itemize}

\CNkeepfig{m23-up-path}{5}

\CNsection{}{Complete 5G User Plane Path Diagram}

\CNfigure{m23-up-path}{Complete 5G user-plane path}

\Needspace{26\baselineskip}

\CNsection{}{Summary Table: User Plane Across Generations}

{\def\LTcaptype{none} % do not increment counter
\Needspace{20\baselineskip}
\begin{longtable}[c]{@{}cccc@{}}
\toprule\noalign{}
\rowcolor{CNink}\CNth{Aspect} & \CNth{3G (UMTS)} & \CNth{4G (LTE)} & \CNth{5G (NR)} \\
\CNheadrulesep
\endhead
\bottomrule\noalign{}
\endlastfoot
Tunnel Protocol & GTP-U & GTP-U & GTP-U \\
RAN\ensuremath{\to}Core Interface & Iu-PS & S1-U & N3 \\
Core Internal & Gn (SGSN\ensuremath{\leftrightarrow}GGSN) & S5/\allowbreak{}S8 (SGW\ensuremath{\leftrightarrow}PGW) & N9 (UPF\ensuremath{\leftrightarrow}UPF) \\
Core\ensuremath{\to}DN & Gi & SGi & N6 \\
QoS Mechanism & Traffic classes & Bearers (QCI) & QoS Flows (QFI) \\
Radio QoS Layer & — & — & SDAP \\
Gateway Function & GGSN & P-GW & UPF \\
Control of UP & — & — & PFCP (N4) \\
Edge Support & No & Limited & Native (UL CL, BP) \\
\end{longtable}
}

\CNboxsection{Key Takeaways}

\begin{CNbox}[unbreakable]{sum}{Key Takeaways}

\begin{enumerate}
\def\labelenumi{\arabic{enumi}.}
\tightlist
\item
  \textbf{GTP-U} is the universal user plane tunnel — same protocol from 3G to 5G
\item
  \textbf{TEID} identifies each tunnel endpoint — assigned by the receiver
\item
  \textbf{5G separates control and user plane} (CUPS) — SMF controls UPF via PFCP/\allowbreak{}N4
\item
  \textbf{SDAP} is the new 5G layer that maps QoS flows to radio bearers
\item
  \textbf{UPF rules (PDR/\allowbreak{}FAR/\allowbreak{}QER/\allowbreak{}URR/\allowbreak{}BAR)} provide programmable packet processing
\item
  \textbf{Edge computing} is natively supported via UL Classifier and Branching Point
\item
  User plane optimization (DC, multi-UPF, edge) enables diverse 5G use cases
\end{enumerate}

\end{CNbox}
